\chapter{Link-Cut Trees and Dynamic Trees}

Heavy-light decomposition (Chapter~19) taught us to answer path queries on a
\emph{fixed} tree in $\Oh(\log^2 n)$ or $\Oh(\log n)$ time. But what if the tree
itself is moving under our feet? Suppose edges are inserted and deleted online,
the forest reshapes between every query, and we still want to ask ``what is the
maximum weight on the path from $u$ to $v$?'' or ``are $u$ and $v$ in the same
tree?'' Heavy-light decomposition collapses the moment the tree changes: the
decomposition into heavy chains is a global property of the whole tree, and a
single edge insertion can shatter it.

The \emph{link-cut tree} of Sleator and Tarjan is the answer. It maintains a
forest of rooted trees under two structural operations,
\[
  \textsc{Link}(u,v) \quad\text{(make $u$ a child of $v$)}
  \qquad\text{and}\qquad
  \textsc{Cut}(u) \quad\text{(detach $u$ from its parent)},
\]
together with path operations (query the path to the root, add a constant along
it, find the root, find the minimum), each in $\Oh(\log n)$ \emph{amortized}
time. It is the data structure that powers Dinic-style blocking-flow algorithms
at their fastest, online dynamic connectivity on forests, and a wide range of
dynamic-tree problems in algorithm design and competitive programming.

This chapter builds the structure from the ground up. We start with the idea of
\emph{preferred paths} -- a self-adjusting cousin of heavy-light chains -- show
how to store each path in a balanced BST (a splay tree), define the single
master operation \textsc{Access} from which everything else follows, and then
prove the $\Oh(\log n)$ amortized bound. We trace a small forest by hand the
whole way through.

\begin{historybox}
The link-cut tree was invented in 1983 by Daniel Sleator and Robert Tarjan at
Bell Labs and Stanford. Their driving motivation was not an abstract tree puzzle,
but maximum network flow. In Dinic's blocking-flow algorithm (1970) and Karzanov's
preflow-push methods (1974), each phase repeatedly pushes flow along augmenting paths
in a layered graph. With standard depth-first search, finding and bottlenecking paths
took $\Oh(V \cdot E)$ per phase, totaling $\Oh(V^2 E)$. Sleator and Tarjan realized
that if one could maintain the flow-carrying paths as a collection of dynamic trees
under edge insertions (\textsc{Link}) and edge deletions (\textsc{Cut}) in $\Oh(\log V)$
time per update, Dinic's overall bound would drop to $\Oh(V E \log V)$!

Every naive attempt had stalled. In the late 1970s, the known balanced tree machinery
(AVL trees, 2-3 trees) balanced binary search trees by local rotations. But if you try
to balance the represented tree directly, a single \textsc{Cut} near the root can orphan
$\Th(n)$ nodes, and local rotations would change the vertex degrees and topology of
the graph being modeled. You cannot rebalance a graph problem without altering the
problem itself!

Sleator and Tarjan's breakthrough was a two-level decoupling: leave the
represented tree's graph topology completely untouched, but decompose its paths into
depth-ordered chains. Each chain is stored inside an auxiliary balanced tree keyed
strictly by depth. To keep the auxiliary trees balanced without complex rebalancing
invariants across splices, they invented \emph{Splay Trees} (published in 1985),
relying on self-adjusting rotations that amortize to $\Oh(\log n)$ per operation.
While biased search trees provided worst-case bounds in their 1983 paper, the splay-based
version became the universal standard for dynamic tree algorithms. An open challenge
remains obtaining worst-case logarithmic bounds with data structures as simple and
cache-efficient as splay-based link-cut trees.
\end{historybox}

\chapterline

\section{What we are maintaining}
\label{ch22:setup}

\begin{focusbox}
We maintain a forest $F$ of rooted trees on $n$ nodes. Each node may carry a
value; we support:
\begin{itemize}
  \item $\textsc{MakeTree}(v)$: create a one-node tree (a new root).
  \item $\textsc{Link}(u,v)$: attach root $u$ as a child of node $v$. Requires
        that $u$ is a tree root and $u,v$ lie in different trees.
  \item $\textsc{Cut}(u)$: remove the edge from $u$ to its parent, splitting one
        tree into two. Requires that $u$ is not a root.
  \item $\textsc{FindRoot}(v)$: return the root of $v$'s tree (so
        $\textsc{FindRoot}(u)=\textsc{FindRoot}(v)$ tests connectivity).
  \item $\textsc{PathAggregate}(v)$: return an aggregate (sum, min, max) of the
        values on the path from $v$ up to its root, and optionally update them.
  \item $\textsc{Evert}(v)$ (re-root): make $v$ the root of its tree, reversing
        the orientation of the path from $v$ to the old root. This upgrades the
        ``path to root'' primitives into ``path between any $u$ and $v$.''
\end{itemize}
\end{focusbox}

These are the operations on the \emph{represented forest} -- the actual forest
the user cares about. The whole trick is that we never store the represented
forest directly with parent pointers, because a single \textsc{Cut} near the top
could orphan $\Th(n)$ descendants and a single path query could walk $\Th(n)$
edges. Instead we store an \emph{auxiliary} structure whose shape only loosely
follows the represented forest, and we re-balance it lazily.

\begin{intuitionbox}
Think of two layers.
\begin{itemize}
  \item The \textbf{represented forest}: the rooted trees the algorithm above us
        reasons about. Nodes here have a true parent.
  \item The \textbf{auxiliary forest}: a collection of balanced binary search
        trees, one per ``preferred path,'' glued together by special
        \emph{path-parent} pointers. This is what we physically store and
        rotate.
\end{itemize}
Every represented-forest operation is translated into a handful of operations on
the auxiliary forest. The key guarantee is that, with splay trees as the
auxiliary BSTs, every translation costs $\Oh(\log n)$ amortized.
\end{intuitionbox}

\section{Preferred paths: a self-adjusting chain decomposition}
\label{ch22:preferred-paths}

Recall heavy-light decomposition. Each node points to exactly one ``heavy''
child (the one with the largest subtree), and the heavy edges link up into chains
that cover the tree. A root-to-node path crosses $\Oh(\log n)$ chains because
each time you leave a chain you at least double the subtree size.

Link-cut trees use the same ``one special child per node'' idea, but the choice
of special child is driven by \emph{access history} rather than subtree size.

\begin{definition}[Preferred child and preferred path]
\label{ch22:def-preferred}
At each node $v$ we designate at most one child as the \emph{preferred child}.
Concretely, the preferred child of $v$ is the child through which the most
recent \textsc{Access} to a descendant of $v$ descended; if the last access to
$v$'s subtree ended at $v$ itself (or no such access has happened), $v$ has no
preferred child. The \emph{preferred edges} are the edges from each node to its
preferred child. They decompose the tree into vertex-disjoint
\emph{preferred paths}.
\end{definition}

So preferred paths are exactly heavy chains in spirit, but they migrate. After
you query the path to node $x$, the entire root-to-$x$ path becomes one
preferred path; everything that used to branch off it loses its preferred-edge
status.

\begin{intuitionbox}
A preferred path is the ``hot'' route your last query took. Storing each hot
route in a balanced BST lets you grab any contiguous chunk of it -- a path query
-- in $\Oh(\log n)$. When the next query takes a different route, you splice the
BSTs accordingly. Self-adjusting access patterns are precisely what splay trees
were designed to exploit, which is why they are the BST of choice here.
\end{intuitionbox}

\subsection{The Two-Layer Architecture: Represented Tree vs. Auxiliary Splay Forest}

Figure~\ref{ch22:fig-two-layers} visualizes this dual-layer reality. In the represented forest (Lane~1), a tree of nine nodes has solid blue edges marking the current preferred path $1\!-\!3\!-\!6\!-\!7$ (the route traversed by the most recent access to node $7$). Dashed gray edges are non-preferred branches. 

In the auxiliary forest (Lane~2), each preferred path lives in its own binary search tree, keyed strictly by depth in the represented forest. Shallower nodes sort to the left; deeper nodes sort to the right. Dashed red arrows show path-parent pointers: each leaves the root of an auxiliary BST and points to the represented-forest parent of the shallowest node on that path.

\begin{figure}[H]
\centering
\begin{tikzpicture}[>=Stealth,
  rnode/.style={circle,draw=sourceblue,fill=sourceblue!12,thick,minimum size=6.5mm,inner sep=0pt,font=\small},
  snode/.style={circle,draw=sourceblue,fill=sourceblue!18,thick,minimum size=6.5mm,inner sep=0pt,font=\small},
  pnode/.style={circle,draw=quotientorange,fill=quotientorange!18,thick,minimum size=6.5mm,inner sep=0pt,font=\small},
  prefedge/.style={very thick,sourceblue},
  nprefedge/.style={thick,dashed,edgegray},
  bstedge/.style={thick,sourceblue,-{Stealth}},
  ppedge/.style={thick,dashed,trapred,-{Stealth}}
]
  % Lane 1: Represented Tree
  \node[font=\footnotesize\bfseries\scshape,text=sourceblue!80!black,anchor=west] at (-4.2,2.8) {Lane 1: Represented Forest (Logical Tree)};
  \draw[gray!30,dashed] (-4.2,2.6) -- (6.0,2.6);

  \node[rnode] (r1) at (0,2.1) {1};
  \node[rnode] (r2) at (-2.2,1.1) {2};
  \node[rnode] (r3) at (1.6,1.1) {3};
  \node[rnode] (r4) at (-2.2,0.0) {4};
  \node[rnode] (r5) at (0.4,0.0) {5};
  \node[rnode] (r6) at (2.8,0.0) {6};
  \node[rnode] (r8) at (1.8,-1.1) {8};
  \node[rnode] (r7) at (3.8,-1.1) {7};
  \node[rnode] (r9) at (3.8,-2.2) {9};

  \draw[prefedge] (r1)--(r3);
  \draw[prefedge] (r3)--(r6);
  \draw[prefedge] (r6)--(r7);
  \draw[nprefedge] (r1)--(r2);
  \draw[nprefedge] (r2)--(r4);
  \draw[nprefedge] (r3)--(r5);
  \draw[nprefedge] (r6)--(r8);
  \draw[nprefedge] (r7)--(r9);

  \node[font=\scriptsize,sourceblue,anchor=west] at (4.3,0.3) {Preferred path: $\{1,3,6,7\}$};
  \node[font=\scriptsize,edgegray,anchor=west] at (4.3,-0.3) {Dashed: non-preferred};

  % Transition divider
  \draw[gray!40,thick] (-4.2,-2.6) -- (6.0,-2.6);

  % Lane 2: Auxiliary Forest
  \node[font=\footnotesize\bfseries\scshape,text=structurepurple!80!black,anchor=west] at (-4.2,-3.0) {Lane 2: Auxiliary Splay Forest (Physical BSTs in Depth Order)};
  \draw[gray!30,dashed] (-4.2,-3.2) -- (6.0,-3.2);

  % Splay tree for {1,3,6,7}, root at 6
  \node[snode] (s6) at (2.4,-3.9) {6};
  \node[snode] (s3) at (1.2,-4.8) {3};
  \node[snode] (s1) at (0.0,-5.7) {1};
  \node[snode] (s7) at (3.6,-4.8) {7};
  \draw[bstedge] (s6)--(s3);
  \draw[bstedge] (s3)--(s1);
  \draw[bstedge] (s6)--(s7);

  % In-order depth ruler for {1,3,6,7}
  \draw[<->,gray!60,thick] (-0.2,-6.3) -- (3.8,-6.3) node[midway,below=1pt,font=\tiny,black] {in-order = depth order: $1 < 3 < 6 < 7$};

  % Splay tree for {2,4}, root at 4
  \node[pnode] (s4) at (-3.1,-4.5) {4};
  \node[pnode] (s2) at (-4.1,-5.5) {2};
  \draw[bstedge] (s4)--(s2);

  % Singletons
  \node[pnode] (s5) at (-1.1,-4.0) {5};
  \node[pnode] (s8) at (0.6,-3.9) {8};
  \node[pnode] (s9) at (4.8,-4.5) {9};

  % Path-parent pointers
  \draw[ppedge] (s4) to[bend right=22] node[midway,below=1pt,font=\scriptsize,text=trapred] {pp} (s1);
  \draw[ppedge] (s5) to[bend right=15] node[midway,below=1pt,font=\scriptsize,text=trapred] {pp} (s3);
  \draw[ppedge] (s8) to[bend left=15] node[midway,above=1pt,font=\scriptsize,text=trapred] {pp} (s6);
  \draw[ppedge] (s9) to[bend right=16] node[midway,above=1pt,font=\scriptsize,text=trapred] {pp} (s7);
\end{tikzpicture}
\caption{The two layers of a Link-Cut Tree. \textbf{Top (Lane~1):} The represented forest where thick blue edges mark the current preferred path $\{1,3,6,7\}$, and dashed gray edges are non-preferred branches. \textbf{Bottom (Lane~2):} The auxiliary forest storing each preferred path as an independent splay tree ordered strictly by depth (shallowest at far left, deepest at far right). Dashed red arrows are unidirectional path-parent pointers leaving auxiliary roots.}
\label{ch22:fig-two-layers}
\end{figure}

The preferred paths here are $\{1,3,6,7\}$, $\{2,4\}$, $\{5\}$, $\{8\}$, $\{9\}$.
Five paths, five auxiliary BSTs. The dashed edges become path-parent pointers:
the BST holding $\{2,4\}$ knows that its topmost node ($2$) hangs off node $1$ in
the represented forest, but node $1$ does \emph{not} point back down to it.

\section{The auxiliary structure: one splay tree per preferred path}
\label{ch22:aux}

We store each preferred path in a binary search tree keyed by \emph{depth in the
represented forest}: shallower nodes are ``smaller.'' So an in-order traversal of
the auxiliary BST visits the path's nodes from the root end to the deep end. The
leftmost node of the BST is the shallowest node of the path; the rightmost is the
deepest.

\begin{formalbox}
\textbf{Invariant (ordering).} For a preferred path $p$, its auxiliary BST stores
exactly the nodes of $p$, ordered by represented-forest depth. The
\emph{symmetric order} (in-order) of the BST equals the top-to-bottom order along
$p$.

\textbf{Invariant (gluing).} The root $r$ of each auxiliary BST stores a
\emph{path-parent pointer} to the represented-forest parent of the shallowest
node on its path (or \texttt{null} if that node is a forest root). No path-parent
pointer points \emph{into} the middle of another BST; it points at a node, and
that node's BST is the parent path. These one-way pointers are the only links
between distinct auxiliary trees.
\end{formalbox}

The two kinds of pointers are easy to confuse, so pin them down:

\begin{itemize}
  \item \textbf{BST child pointers} (\texttt{left}, \texttt{right}) live
        \emph{inside} one auxiliary tree and obey the depth ordering. They are
        rotated freely by splay operations and carry no represented-forest
        meaning beyond the ordering invariant.
  \item \textbf{Path-parent pointers} live \emph{between} auxiliary trees. They
        encode the non-preferred edges of the represented forest.
\end{itemize}

\begin{trapbox}
A path-parent pointer is stored at the \emph{root} of an auxiliary tree and is
unidirectional: child path knows its parent path, but the parent stores no list
of its hanging children. This asymmetry is what keeps \textsc{Cut} cheap -- you
never have to scan or update a parent's child list -- but it also means
``which preferred paths hang off me?'' is not a question the structure answers
directly. You only ever walk path-parent pointers \emph{upward}.
\end{trapbox}

The dashed path-parent arrows are the only links between distinct trees, and they
run strictly upward toward the represented root. Walking them upward is exactly
what \textsc{Access} does.

\subsection{Why a BST keyed by depth}

Because the key is depth, a contiguous prefix of the path (everything shallower
than some node $v$) is exactly the set of BST nodes less than $v$ -- the left
subtree after we bring $v$ to the root. That is what makes path-to-root
aggregates $\Oh(\log n)$: splay $v$ to the top of its auxiliary tree, and its
left subtree is precisely the part of the current preferred path that lies above
$v$.

\section{Splay trees, briefly}
\label{ch22:splay}

The auxiliary BSTs are \emph{splay trees}: self-adjusting binary search trees
that move an accessed node to the root by a sequence of rotations called
\emph{splaying}. We will not re-derive the full splay analysis here (the
companion intro book covers basic BSTs), but we need the operation and its
amortized guarantee.

A single rotation at $x$ (with parent $p$) is the usual BST rotation; splaying
repeats double rotations until $x$ is the root:

\begin{center}
\begin{tikzpicture}[
  tnode/.style={circle,draw=nodeblue,thick,fill=blue!6,minimum size=6.2mm,inner sep=0pt,font=\footnotesize},
  edge from parent/.style={draw=nodeblue,thick},
  level distance=9.5mm, sibling distance=12mm]
  \begin{scope}
    \node[tnode] (p) {$p$}
      child {node[tnode] (x) {$x$}
        child {node[tnode] (a) {$A$}}
        child {node[tnode] (b) {$B$}}}
      child {node[tnode] (c) {$C$}};
    \node[font=\small,text=nodeblue] at (-0.3,-2.45) {zig at $x$};
  \end{scope}
  \begin{scope}[xshift=58mm]
    \node[tnode] (x2) {$x$}
      child {node[tnode] (a2) {$A$}}
      child {node[tnode] (p2) {$p$}
        child {node[tnode] (b2) {$B$}}
        child {node[tnode] (c2) {$C$}}};
    \node[font=\small,text=nodeblue] at (0.3,-2.45) {result};
  \end{scope}
  \draw[-{Stealth},very thick,nodeblue] (2.4,-1.0) -- (3.8,-1.0);
\end{tikzpicture}
\end{center}

\begin{formalbox}
\textbf{Splay amortized bound.} Assign each node $x$ a weight $w(x)>0$ and define
its \emph{size} $s(x)=\sum_{y\in\text{subtree}(x)} w(y)$ and \emph{rank}
$r(x)=\log_2 s(x)$. The potential of a splay tree is $\Phi=\sum_x r(x)$.
\emph{Access Lemma:} splaying a node $x$ in a tree with root $t$ costs at most
$3\bigl(r(t)-r(x)\bigr)+1$ amortized rotations. Consequently a sequence of $m$
operations on a tree of $n$ nodes runs in $\Oh\bigl((m+n)\log n\bigr)$.
\end{formalbox}

The weight flexibility in the Access Lemma is exactly the hook we will use:
choosing $w(x)$ to count the \emph{whole represented subtree} of $x$, not just
its auxiliary subtree, is what turns the per-splay bound into a per-\textsc{Access}
bound across all the auxiliary trees at once.

\section{Access: the master operation}
\label{ch22:access}

Everything reduces to one operation. $\textsc{Access}(v)$ does two things:

\begin{enumerate}
  \item It makes the path from $v$ up to the root of its represented tree the
        \emph{preferred path}, and makes $v$ the deepest node on a preferred path
        that starts at the represented root.
  \item It splays $v$ to the root of its auxiliary tree, so that afterwards the
        single auxiliary tree containing the represented root has $v$ at its top,
        and $v$'s left subtree is exactly the represented path from the root down
        to $v$.
\end{enumerate}

After $\textsc{Access}(v)$, the node $v$ has no path-parent pointer (it is on the
path containing the represented root), and the auxiliary tree rooted at $v$
\emph{is} the represented root-to-$v$ path stored in depth order.

\subsection{How access walks up}

We splay $v$ within its own auxiliary tree, then repeatedly hop up a path-parent
pointer, splay the parent, and splice. Each hop turns one non-preferred edge into
a preferred edge (and the old preferred child of the parent into a
non-preferred one). The splicing is pure BST surgery: at the parent $w$ we splay
$w$, detach its current right subtree (the part of $w$'s path \emph{deeper} than
$w$, which is losing preferred status -- it becomes a separate auxiliary tree
with a path-parent pointer back to $w$), and attach the tree we just came from as
$w$'s new right subtree.

\begin{algorithm}[H]
\caption{$\textsc{Access}(v)$ -- make $v$ the deep end of the root's preferred path}
\label{ch22:alg-access}
\DontPrintSemicolon
\SetKwFunction{Splay}{Splay}
\SetKwData{pp}{pathparent}
$\Splay(v)$\;
\If{$v.\mathrm{right}\neq\textsc{null}$}{
  \tcp{detach $v$'s path below $v$: it becomes its own preferred path}
  $v.\mathrm{right}.\pp \leftarrow v$\;
  $v.\mathrm{right}.\mathrm{parent} \leftarrow \textsc{null}$\;
  $v.\mathrm{right} \leftarrow \textsc{null}$\;
}
recompute aggregate at $v$\;
\While{$v.\pp \neq \textsc{null}$}{
  $w \leftarrow v.\pp$\;
  $\Splay(w)$\;
  \If{$w.\mathrm{right}\neq\textsc{null}$}{
    \tcp{$w$'s old preferred child below it loses preferred status}
    $w.\mathrm{right}.\pp \leftarrow w$\;
    $w.\mathrm{right}.\mathrm{parent} \leftarrow \textsc{null}$\;
  }
  $w.\mathrm{right} \leftarrow v$\;
  $v.\mathrm{parent} \leftarrow w$\;
  $v.\pp \leftarrow \textsc{null}$\;
  recompute aggregate at $w$\;
  $\Splay(v)$ \tcp*{bring $v$ back to the top}
}
\Return $v$\;
\end{algorithm}

\begin{whybox}
Why does this leave the invariants intact? Each iteration of the loop dissolves
one non-preferred edge (the path-parent from $v$'s path to $w$) and creates one
(from $w$'s old right child back to $w$). The depth ordering is preserved because
$w$ is shallower than every node of $v$'s former tree (it is $v$'s
represented-forest ancestor through the spliced edge), so attaching $v$'s tree as
$w$'s \emph{right} subtree keeps in-order $=$ top-to-bottom. When the loop ends,
$v$ has no path-parent, meaning its auxiliary tree contains the represented root;
and after the final \textsc{Splay}, $v$ sits at that tree's root with the whole
root-to-$v$ path in its left subtree.
\end{whybox}

\begin{recoverybox}
\textbf{Pause and think: what does $\textsc{Access}(v)$ accomplish in the represented tree?}
In the represented tree, $\textsc{Access}(v)$ makes every edge along the path from the root down to $v$ preferred, while making the edge leaving $v$ toward any deeper child non-preferred. In the auxiliary forest, this merges all the depth-ordered splay trees along the root-to-$v$ path into a single splay tree, places $v$ at the root of that tree, and ensures that $v$'s right child is empty (no deeper nodes on the preferred path). Everything on the root-to-$v$ path now lives in $v$ and its left subtree.
\end{recoverybox}

\subsection{A hand trace}

\begin{labbox}
Start from the nine-node forest of Section~\ref{ch22:preferred-paths}, with
preferred path $1\!-\!3\!-\!6\!-\!7$. The auxiliary trees, drawn as their in-order
(equivalently, depth order):

\begin{center}
\begin{tabular}{ll}
\toprule
preferred path & auxiliary BST and its path-parent\\
\midrule
$\{1,3,6,7\}$ & BST over $1<3<6<7$; path-parent $=\textsc{null}$ (holds the root)\\
$\{2,4\}$     & BST over $2<4$; path-parent of top node $2$ is $1$\\
$\{5\}$       & singleton; path-parent of $5$ is $3$\\
$\{8\}$       & singleton; path-parent of $8$ is $6$\\
$\{9\}$       & singleton; path-parent of $9$ is $7$\\
\bottomrule
\end{tabular}
\end{center}

Now call $\textsc{Access}(4)$. Walk it:
\begin{enumerate}
  \item Splay $4$ inside its tree $\{2,4\}$; now $4$ is the BST root with left
        child $2$. Node $4$ has no right child, so nothing to detach.
  \item The path-parent now sits at the root $4$ and equals $1$ (it followed the
        topmost path node $2$'s old edge to $1$). Hop to $w=1$.
  \item Splay $1$ in the big tree $\{1,3,6,7\}$. Now $1$ is the BST root; its
        right subtree is $\{3,6,7\}$ (everything deeper than $1$).
  \item $1$ has a right child, so the path $\{3,6,7\}$ loses preferred status:
        set its root's path-parent to $1$, detach it. Attach $4$'s tree as $1$'s
        new right child. The preferred path containing the root is now
        $1\!-\!2\!-\!4$.
  \item $4$ has no further path-parent. Splay $4$ to the top. Done.
\end{enumerate}
Two path-parent rewires, one access. The resulting preferred edges:
\end{labbox}

\begin{figure}[H]
\centering
\begin{tikzpicture}[>=Stealth,
  rnode/.style={circle,draw=sourceblue,fill=sourceblue!12,thick,minimum size=6mm,inner sep=0pt,font=\scriptsize},
  snode/.style={circle,draw=quotientorange,fill=quotientorange!18,thick,minimum size=6mm,inner sep=0pt,font=\scriptsize},
  pref/.style={very thick,sourceblue},
  npref/.style={dashed,edgegray,thick},
  ppedge/.style={thick,dashed,trapred,-{Stealth}}
]
  % Panel 1: Before
  \begin{scope}[xshift=0cm]
    \node[font=\footnotesize\bfseries\scshape,text=sourceblue!80!black] at (0.6,0.6) {Step 1: Before $\textsc{Access}(4)$};
    \node[rnode] (a1) at (0,0) {1};
    \node[snode] (a2) at (-1.5,-0.9) {2};
    \node[rnode] (a3) at (1.2,-0.9) {3};
    \node[snode] (a4) at (-1.5,-1.9) {4};
    \node[rnode] (a5) at (0.3,-1.9) {5};
    \node[rnode] (a6) at (2.1,-1.9) {6};
    \node[rnode] (a8) at (1.4,-2.8) {8};
    \node[rnode] (a7) at (2.8,-2.8) {7};
    \node[rnode] (a9) at (2.8,-3.7) {9};
    \draw[pref] (a1)--(a3); \draw[pref] (a3)--(a6); \draw[pref] (a6)--(a7);
    \draw[npref] (a1)--(a2); \draw[npref] (a2)--(a4);
    \draw[npref] (a3)--(a5); \draw[npref] (a6)--(a8); \draw[npref] (a7)--(a9);
    \node[font=\scriptsize,text=sourceblue,anchor=north] at (0.6,-4.0) {Preferred: $\{1,3,6,7\}$};
  \end{scope}

  % Arrow between panels
  \draw[->,very thick,structurepurple] (3.5,-1.8) -- (5.0,-1.8) node[midway,above=2pt,font=\scriptsize\bfseries] {splice};

  % Panel 2: After
  \begin{scope}[xshift=7.4cm]
    \node[font=\footnotesize\bfseries\scshape,text=targetgreen!80!black] at (0.6,0.6) {Step 2: After $\textsc{Access}(4)$};
    \node[rnode] (b1) at (0,0) {1};
    \node[rnode] (b2) at (-1.5,-0.9) {2};
    \node[snode] (b3) at (1.2,-0.9) {3};
    \node[rnode] (b4) at (-1.5,-1.9) {4};
    \node[snode] (b5) at (0.3,-1.9) {5};
    \node[snode] (b6) at (2.1,-1.9) {6};
    \node[snode] (b8) at (1.4,-2.8) {8};
    \node[snode] (b7) at (2.8,-2.8) {7};
    \node[snode] (b9) at (2.8,-3.7) {9};
    \draw[pref] (b1)--(b2); \draw[pref] (b2)--(b4);
    \draw[npref] (b1)--(b3);
    \draw[npref] (b3)--(b5); \draw[npref] (b3)--(b6);
    \draw[npref] (b6)--(b8); \draw[npref] (b6)--(b7); \draw[npref] (b7)--(b9);
    \node[font=\scriptsize,text=targetgreen!80!black,anchor=north] at (0.6,-4.0) {Preferred: $\{1,2,4\}$};
  \end{scope}
\end{tikzpicture}
\caption{The self-adjusting transition performed by $\textsc{Access}(4)$. \textbf{Left:} Before access, the root path $1\!-\!3\!-\!6\!-\!7$ is preferred. \textbf{Right:} Splaying $4$ and ascending path-parent pointers to $1$ swaps the preferred child of $1$ from $3$ to $2$. The old chain $\{3,6,7\}$ is severed into a separate auxiliary tree hanging off $1$ via a path-parent pointer, and $1\!-\!2\!-\!4$ becomes the active preferred path.}
\label{ch22:fig-access-transition}
\end{figure}

\section{Everything else in terms of access}
\label{ch22:derived}

Once \textsc{Access} is in hand, the represented-forest operations are short.

\paragraph{FindRoot.}
$\textsc{Access}(v)$ puts $v$'s root path into one auxiliary tree with $v$ at the
top. The represented root is the shallowest node, i.e.\ the leftmost node of that
tree. Walk left children to the bottom, then splay it (to keep the tree balanced
and the amortized bound intact) and return it.

\begin{algorithm}[H]
\caption{$\textsc{FindRoot}(v)$}
\label{ch22:alg-findroot}
\DontPrintSemicolon
\SetKwFunction{Access}{Access}\SetKwFunction{Splay}{Splay}
\Access{$v$}\;
$r \leftarrow v$\;
\While{$r.\mathrm{left}\neq\textsc{null}$}{ $\textsc{push}(r)$;\ \ $r \leftarrow r.\mathrm{left}$ \tcp*{descend to the shallowest node} }
$\Splay(r)$\;
\Return $r$\;
\end{algorithm}

\paragraph{Link.}
$\textsc{Link}(u,v)$ with $u$ a represented root: in the re-rooting variant,
\textsc{Evert} $u$ so it becomes the root of its own tree, access $v$, then set
$u$'s path-parent to $v$, hanging $u$'s whole tree off $v$ as a non-preferred
child.

\begin{algorithm}[H]
\caption{$\textsc{Link}(u,v)$ -- attach $u$ under $v$ (re-rooting variant)}
\label{ch22:alg-link}
\DontPrintSemicolon
\SetKwFunction{Access}{Access}\SetKwFunction{Evert}{Evert}
\Evert{$u$} \tcp*{make $u$ the root of its own tree}
\Access{$v$}\;
$u.\mathrm{pathparent} \leftarrow v$ \tcp*{$u$ hangs off $v$ as a non-preferred child}
\end{algorithm}

\paragraph{Cut.}
$\textsc{Cut}(u)$ removes the edge to $u$'s parent. Access $u$; its left subtree
is exactly the represented path strictly above $u$, i.e.\ $u$'s ancestors. Detach
that left subtree -- it becomes its own tree, rooted at the old represented root
-- and $u$ becomes a new represented root.

\begin{algorithm}[H]
\caption{$\textsc{Cut}(u)$ -- detach $u$ from its parent}
\label{ch22:alg-cut}
\DontPrintSemicolon
\SetKwFunction{Access}{Access}
\Access{$u$}\;
\If{$u.\mathrm{left}\neq\textsc{null}$}{
  $u.\mathrm{left}.\mathrm{parent} \leftarrow \textsc{null}$\;
  $u.\mathrm{left} \leftarrow \textsc{null}$\;
}
recompute aggregate at $u$\;
\end{algorithm}

\paragraph{Evert (re-root).}
To make $v$ the root, access $v$ (bringing the root-to-$v$ path into one tree with
$v$ at the top and the path in its left subtree), then \emph{reverse} that path.
Reversal is the standard lazy ``flip'' bit on the auxiliary tree root that swaps
left and right subtrees on push-down -- the same trick used in a treap or splay
tree to support sequence reversal. After flipping at $v$, depth order is inverted,
so $v$ (formerly deepest) becomes shallowest, i.e.\ the new root.

\begin{algorithm}[H]
\caption{$\textsc{Evert}(v)$ -- make $v$ the root of its tree}
\label{ch22:alg-evert}
\DontPrintSemicolon
\SetKwFunction{Access}{Access}
\Access{$v$}\;
$v.\mathrm{flip} \leftarrow \lnot\, v.\mathrm{flip}$ \tcp*{lazily reverse the root-to-$v$ path}
\end{algorithm}

\paragraph{Path queries.}
For a query on the path between $u$ and $v$ (sum, min, max, or a range update):
$\textsc{Evert}(u)$ then $\textsc{Access}(v)$. Now the auxiliary tree rooted at
$v$ holds \emph{exactly} the nodes of the $u$--$v$ path, in order. Read the
aggregate stored at the root $v$; for an update, set a lazy tag at $v$. If you
only need root-to-$v$ aggregates (rooted variant, no re-rooting), a bare
$\textsc{Access}(v)$ already isolates the path in $v$'s tree.
Figure~\ref{ch22:fig-cut-and-evert} shows how \textsc{Cut}$(u)$ and
$\textsc{Evert}(u)+\textsc{Access}(v)$ act directly on the auxiliary splay tree.

\begin{figure}[H]
\centering
\begin{tikzpicture}[>=Stealth,
  snode/.style={circle,draw=sourceblue,fill=sourceblue!15,thick,minimum size=6.2mm,inner sep=0pt,font=\scriptsize},
  rnode/.style={snode,draw=structurepurple,very thick,fill=purple!15}
]
% Panel (a): Cut(u)
\begin{scope}[xshift=0cm]
  \node[font=\footnotesize\bfseries,text=sourceblue!85!black] at (1.8, 2.5) {(a) $\textsc{Cut}(u)$ after $\textsc{Access}(u)$};
  \node[rnode] (u) at (2.2, 1.5) {$u$};
  \node[snode] (p) at (1.0, 0.4) {$p$};
  \node[snode] (r) at (0.2, -0.6) {$r$};
  \draw[thick,sourceblue,->] (p) -- (r);
  \draw[thick,trapred,dashed,->] (u) -- node[above left=-1pt,font=\scriptsize,text=trapred]{sever \texttt{u.left}} (p);
  \draw[draw=sourceblue!60,dashed,rounded corners=3pt] (-0.35, -1.1) rectangle (1.55, 0.9);
  \node[font=\scriptsize,text=sourceblue!80!black] at (0.6, -1.4) {ancestors ($<\mathrm{dep}(u)$)};
  \node[font=\scriptsize,text=edgegray,right=3pt of u] {$\texttt{u.right}=\textsc{null}$};
\end{scope}

% Panel (b): Evert(u) + Access(v)
\begin{scope}[xshift=6.8cm]
  \node[font=\footnotesize\bfseries,text=targetgreen!80!black] at (2.0, 2.5) {(b) $\textsc{Evert}(u)$ then $\textsc{Access}(v)$};
  \node[rnode,draw=targetgreen!80!black,fill=targetgreen!18] (v) at (2.2, 1.5) {$v$};
  \node[snode] (w) at (1.1, 0.5) {$w$};
  \node[snode] (u2) at (0.2, -0.5) {$u$};
  \node[snode] (z) at (1.9, -0.5) {$z$};
  \draw[thick,targetgreen!75!black,->] (v) -- (w);
  \draw[thick,targetgreen!75!black,->] (w) -- (u2);
  \draw[thick,targetgreen!75!black,->] (w) -- (z);
  \draw[draw=targetgreen!70!black,dashed,rounded corners=3pt] (-0.35, -1.1) rectangle (2.75, 2.0);
  \node[font=\scriptsize\bfseries,text=targetgreen!75!black] at (1.2, -1.4) {auxiliary BST at $v$ = exact $u\!\leftrightarrow\!v$ path};
\end{scope}
\end{tikzpicture}
\caption{Auxiliary-tree surgery for \textsc{Cut} and path queries.
(a)~After $\textsc{Access}(u)$, $u$ sits at the auxiliary root with $\texttt{u.right}=\textsc{null}$
and all ancestors on the root-to-$p$ path in $\texttt{u.left}$; nulling $\texttt{u.left}$ cuts edge $(p,u)$.
(b)~$\textsc{Evert}(u)$ makes $u$ the represented root (flipping its splay tree), after which
$\textsc{Access}(v)$ gathers the entire $u\!\leftrightarrow\!v$ path into the single auxiliary BST rooted at $v$.}
\label{ch22:fig-cut-and-evert}
\end{figure}

\begin{trapbox}
\textsc{Evert} only works if your aggregate is \emph{reversible}. Sum, min, and
max are symmetric, so flipping the path order does not change them. But a
\emph{non-commutative} aggregate (matrix products, or ``leftmost occurrence'')
must store \emph{two} accumulators per subtree -- one for each direction -- and
swap them when the flip bit is pushed, exactly as with a reversible segment tree.
Forgetting this is the classic link-cut bug: queries silently return the answer
for the reversed path.
\end{trapbox}

\begin{whybox}
Why is $\textsc{Evert}(u);\textsc{Access}(v)$ exactly the $u$--$v$ path? After
\textsc{Evert}, $u$ is the root of the represented tree, so the represented path
from the root to $v$ \emph{is} the path from $u$ to $v$. \textsc{Access}$(v)$
then gathers precisely the root-to-$v$ path into one auxiliary tree with $v$ on
top. The aggregate at $v$ summarizes its whole auxiliary subtree, which is that
path and nothing else.
\end{whybox}

\begin{recoverybox}
\textbf{Pause and think: why does swapping left and right subtrees reverse the path?}
In an auxiliary splay tree, the binary search tree key is the depth in the represented tree: in-order traversal lists nodes from shallowest to deepest. Flipping the left and right children of every node inverts the in-order sequence completely. The deepest node ($v$) becomes the shallowest, and parent-child orientations along the entire path invert without touching any off-path nodes!
\end{recoverybox}

\section{Correctness: the invariants never break}
\label{ch22:correctness}

\begin{theorem}[Structural correctness]
\label{ch22:thm-correct}
Starting from singletons, any sequence of $\textsc{MakeTree}$, $\textsc{Link}$,
$\textsc{Cut}$, $\textsc{Evert}$, $\textsc{Access}$ (with legal preconditions)
maintains the two invariants of Section~\ref{ch22:aux}, and the represented
forest defined by ``preferred edges $+$ path-parent edges'' equals the abstract
forest these operations specify.
\end{theorem}

\begin{proof}[Proof sketch]
We argue each primitive preserves both invariants; \textsc{Access} is the only
nontrivial one and the others are built on it.

\emph{Splay preserves invariants.} Rotations inside an auxiliary tree change BST
shape but not the set of nodes nor their symmetric order, so the depth-ordering
invariant holds. The path-parent of a tree must always sit at the root; whenever
splaying moves a node to the root we transfer the path-parent pointer with it (a
constant-time fix in the \textsc{Splay} wrapper). The gluing invariant is
otherwise untouched: no path-parent target moves between trees.

\emph{Access preserves invariants.} By the loop-invariant argument in
Section~\ref{ch22:access}: each iteration converts one non-preferred edge into a preferred edge and
one preferred edge into a non-preferred one, with the path-parent pointer
relocated to the new tree root. Attaching the incoming tree as the parent's
\emph{right} subtree respects depth order because the parent is strictly
shallower than every incoming node. The set of represented edges (preferred
$\cup$ path-parent) is unchanged throughout -- we only relabel which are
preferred -- so the represented forest is invariant.

\emph{Link, Cut, Evert} are compositions of \textsc{Access} with one $\Oh(1)$
pointer edit (set a path-parent, null a left child, toggle a flip bit), each of
which matches the abstract semantics: \textsc{Link} adds the single edge
$u\!\to\!v$ as a path-parent; \textsc{Cut} deletes the single edge above $u$ by
nulling the left subtree that holds $u$'s ancestors; \textsc{Evert} reverses the
root-to-$v$ chain, which is precisely re-rooting at $v$. Lazy flip bits are
pushed down before any structural read, so reads see the correct orientation.
\end{proof}

\begin{remark}
The lazy flip bit deserves the same discipline as in a balanced-BST sequence
structure: \emph{push down} (apply and clear the tag, swapping children, and
forwarding the tag to both children) before you recurse into or rotate around a
node. Splay's rotations must push down the two nodes involved first, or the
ordering invariant silently corrupts.
\end{remark}

\section{The amortized analysis}
\label{ch22:analysis}

Now the payoff: each \textsc{Access} costs $\Oh(\log n)$ amortized, so every
represented-forest operation does too.

The cost of one \textsc{Access} has two parts: (a) the number of splay steps,
dominated by the number of \emph{preferred-child changes} (one splice per loop
iteration), plus (b) the splay work itself. We bound both.

\subsection{Counting preferred-child changes}

\begin{lemma}[Few preferred changes, amortized]
\label{ch22:lem-prefchanges}
Over any sequence of $m$ accesses on an $n$-node forest built by links and cuts,
the total number of preferred-child changes is $\Oh\bigl((m+n)\log n\bigr)$.
\end{lemma}

The slick way to see this is a \emph{heavy-light} counting argument. We give the
intuition; the full charge is folded into the splay potential below.

\begin{intuitionbox}
Classify each represented edge as \emph{heavy} (the child's subtree is more than
half the parent's subtree) or \emph{light} otherwise, exactly as in heavy-light
decomposition. Any root-to-node path crosses $\Oh(\log n)$ light edges, because
each light edge at least halves the subtree size as you go up. Charge each
preferred-child change either to the splay potential (if it happens along a heavy
edge -- and there are few heavy edges on any path) or to a light edge it crosses.
Summing over the path gives $\Oh(\log n)$ changes per access plus a global term
for heavy edges, which telescopes to $\Oh((m+n)\log n)$ overall.
\end{intuitionbox}

\subsection{The splay potential does the heavy lifting}

Here is the clean, self-contained version. Weight every node by $1$, and define
the \emph{size} $s(x)$ of an auxiliary-tree node $x$ to be the total number of
\emph{represented} descendants reachable through $x$ -- including those in
auxiliary subtrees hanging off $x$ by path-parent pointers. With this choice
$s(\text{represented root})=n$. Put $r(x)=\log_2 s(x)$ and $\Phi=\sum_x r(x)$,
summed over all nodes of all auxiliary trees.

\begin{theorem}[Sleator--Tarjan]
\label{ch22:thm-amortized}
With this potential, one \textsc{Access} on an $n$-node link-cut forest costs
$\Oh(\log n)$ amortized. Hence \textsc{FindRoot}, \textsc{Link}, \textsc{Cut},
\textsc{Evert}, and a path query each run in $\Oh(\log n)$ amortized time, and a
sequence of $m$ operations runs in $\Oh\bigl((m+n)\log n\bigr)$ total.
\end{theorem}

\begin{proof}[Proof idea]
Fix the access to $v$. The work is a sum of splays, one per preferred path on the
root-to-$v$ route. Apply the Access Lemma to each splay: the splay of $v$ inside a
path whose auxiliary root is $t$ costs at most $3\bigl(r(t)-r(v)\bigr)+1$
amortized. When we then splice $v$'s tree under the parent $w$ and splay $w$, the
rank $r(w)$ is at least the old $r(t)$ (it is shallower, hence has a larger
represented subtree), so consecutive $3(r(\cdot)-r(\cdot))$ terms
\emph{telescope}:
\[
  \sum_{\text{paths } t_1,t_2,\dots \text{ up to the root}}
    3\bigl(r(t_i)-r(v)\bigr)
  \;\le\; 3\bigl(r(\text{represented root}) - r(v)\bigr) \;+\; (\text{\#splices}).
\]
Since $s(\text{represented root})=n$, we get $r(\text{root})\le\log_2 n$, so the
telescoped sum is $\Oh(\log n)$ plus one unit per splice. The number of splices is
the number of preferred-child changes on this access; the $+1$ units are absorbed
by Lemma~\ref{ch22:lem-prefchanges}, whose total over all $m$ accesses is
$\Oh\bigl((m+n)\log n\bigr)$. Dividing by $m$ gives $\Oh(\log n)$ amortized per
access. Each derived operation is one access plus $\Oh(1)$ extra splay or pointer
work, so the bound carries over.

The $+n$ term accounts for the initial potential and for \textsc{Link}s, which
can raise a represented subtree (and hence ranks) by an amount bounded by the
$\log n$ rank ceiling.
\end{proof}

\begin{whybox}
The single most important line is the telescoping. Because depth order makes each
spliced-in parent shallower, its rank only goes \emph{up} as we climb, so the
$-r(v)$ from one path cancels against the $+r(t)$ of the next. The cost of an
entire access collapses to ``rank of the root minus rank of $v$,'' a quantity
bounded by $\log_2 n$. This is the same engine that makes a single splay
$\Oh(\log n)$; link-cut trees just run it once per preferred path and let the
ranks chain together.
\end{whybox}

\begin{recoverybox}
\textbf{Pause and think: why does access pay for both preferred and non-preferred splices?}
Each step of $\textsc{Access}$ that flips a non-preferred edge into a preferred edge costs real work, but decreases the number of heavy non-preferred edges along the path. By defining the potential as $\sum \log \mathrm{size}(x)$, the telescoping sum of splay operations absorbs the cost of all preferred updates, proving that the amortized cost per operation remains $\Oh(\log n)$ regardless of tree shape!
\end{recoverybox}

\begin{complexitybox}
\textbf{Per operation (amortized):}
\[
\begin{array}{ll}
\textsc{Access},\ \textsc{FindRoot} & \Oh(\log n)\\[2pt]
\textsc{Link},\ \textsc{Cut},\ \textsc{Evert} & \Oh(\log n)\\[2pt]
\text{path aggregate / path update} & \Oh(\log n)\\
\end{array}
\]
\textbf{Space:} $\Oh(n)$ -- a constant number of pointers, an aggregate, and a
lazy tag per node. \textbf{Worst case for a single op:} $\Th(n)$ (a splay can be
linear once), but the amortized average is $\Oh(\log n)$. Variants built on
globally biased or balanced BSTs achieve $\Oh(\log n)$ \emph{worst case} per
operation at the cost of much heavier bookkeeping.
\end{complexitybox}

\paragraph{Theoretical milestones and variants.}
Daniel Sleator and Robert Tarjan introduced link-cut trees in ``A Data Structure
for Dynamic Trees'' (JCSS, 1983) to speed up the blocking-flow phases of
maximum-flow algorithms. The companion paper that gave us splay trees
(Sleator--Tarjan, 1985) supplied the self-adjusting auxiliary structure that
makes the modern link-cut tree so concise to implement. The original 1983 version
used \emph{biased} balanced binary trees for a worst-case bound; the splay-based
version is the one almost everyone codes today. Tarjan and Werneck's experimental
study (2009) compares link-cut trees against Euler-tour trees and top trees and
is the standard reference for practical dynamic-tree performance.

\section{What link-cut trees unlock}
\label{ch22:applications}

\begin{applicationbox}[Blocking flows in $\Oh(VE\log V)$]
Dinic's algorithm finds a blocking flow in each phase by repeatedly pushing along
augmenting paths in the level graph. The naive blocking-flow inner loop is
$\Oh(VE)$ per phase. Sleator and Tarjan's original motivation: maintain the
current set of admissible edges as a link-cut forest where each node points to the
next node on a potential augmenting path, store residual capacities as edge
values, and use \textsc{PathAggregate} to find the bottleneck (minimum residual)
along the whole path at once, then a path update to subtract the flow, then a
\textsc{Cut} on every saturated edge. This drops the blocking-flow phase to
$\Oh(E\log V)$ and the whole max-flow to $\Oh(VE\log V)$, beating plain $\Oh(V^2E)$
Dinic on dense graphs. See Chapter~12.
\end{applicationbox}

\begin{applicationbox}[Online dynamic connectivity on forests]
For a forest under edge insert/delete with connectivity queries,
$\textsc{FindRoot}(u)=\textsc{FindRoot}(v)$ answers ``same tree?'' in
$\Oh(\log n)$. With \textsc{Evert} you also get $\Oh(\log n)$ tree-path queries
(distance, max edge, sum). This is the forest case of the general dynamic
connectivity problem; for \emph{general} graphs one layers a spanning forest of
link-cut (or Euler-tour) trees underneath, replacing tree edges with non-tree
edges on deletion (Holm--de Lichtenberg--Thorup), which Chapter~39
develops further.
\end{applicationbox}

\begin{applicationbox}[Online LCA, MST verification, and more]
Link-cut trees give online \emph{lowest common ancestor}: $\textsc{Access}(u)$
then $\textsc{Access}(v)$ -- the last node at which the second access performed a
splice is the LCA. They support online \emph{maximum-edge-on-path} queries, which
is the inner loop of incremental minimum-spanning-tree maintenance and of MST
verification (is edge $e$ heavier than every tree edge on the cycle it would
close?). Whenever a problem is ``trees that change, plus path questions,''
link-cut trees are the standard tool.
\end{applicationbox}

\begin{connectionbox}
The line from heavy-light decomposition (Chapter~19) to link-cut trees is a single
conceptual move: replace the \emph{static} heavy-chain decomposition with a
\emph{self-adjusting} preferred-path decomposition, and replace the per-chain
segment tree with a per-path splay tree. Euler-tour trees (Chapter~21) solve a
\emph{subset} of the same problems (subtree aggregates, connectivity) but cannot
do arbitrary path aggregates -- the two structures are complementary, and
\emph{top trees} (a generalization built on top of link-cut trees) unify them.

The amortized engine here -- a potential equal to a sum of logs of subtree sizes,
telescoping along a root path -- is the same machinery behind splay trees and,
spiritually, behind the union-by-rank/path-compression analysis of disjoint-set
forests. When you reach Chapters~31--35 and~41, note the contrast: link-cut trees
are a \emph{combinatorial} dynamic-tree structure, whereas dynamic
\emph{electrical} quantities (effective resistances on a changing graph) need the
Laplacian machinery, where ``path'' is replaced by ``current'' and ``min on a
path'' by ``potential drop.''
\end{connectionbox}

\section{Implementation checklist}
\label{ch22:impl}

\begin{focusbox}
A correct link-cut tree node carries:
\begin{itemize}
  \item \texttt{parent}, \texttt{left}, \texttt{right} -- auxiliary-tree links.
        Whether a node is a splay root is detected by ``my parent's children do
        not include me,'' so \texttt{parent} doubles as the path-parent at a
        root.
  \item \texttt{value} and an \texttt{aggregate} over the auxiliary subtree
        (sum/min/max), recomputed on every structural change via \textsc{pull}.
  \item a \texttt{flip} lazy bit (for \textsc{Evert}) and any range-update lazy
        tags, applied via \textsc{push} before recursing or rotating.
\end{itemize}
The two functions you must call religiously: \textsc{push}$(x)$ (apply pending
lazy tags, propagate to children) \emph{before} reading children, and
\textsc{pull}$(x)$ (recompute the aggregate from children) \emph{after} changing
them. Splay rotations must \textsc{push} both nodes first and \textsc{pull} both
after, bottom node before top.
\end{focusbox}

\begin{trapbox}
Three recurring bugs:
\begin{enumerate}
  \item \textbf{Forgetting the root test.} Many implementations fold the
        path-parent into \texttt{parent} and decide ``is this node a splay root''
        by checking whether the parent claims it as a child. Get this predicate
        wrong and splay either stops early (leaving the wrong node at the top) or
        rotates across a path-parent edge (corrupting the gluing invariant).
  \item \textbf{Pushing in the wrong order during rotation.} You must \textsc{push}
        grandparent, parent, then child top-down before rotating, or a stale
        \texttt{flip} bit reverses a subtree you have already read.
  \item \textbf{Non-reversible aggregates under \textsc{Evert}.} As warned above,
        store both directions for non-commutative monoids.
\end{enumerate}
\end{trapbox}

\begin{connectionbox}
\textbf{What This Unlocked:}
Link-cut trees delivered the first practical $\Oh(\log n)$ amortized dynamic forest
data structure. Beyond speeding up maximum network flow to $\Oh(V E \log V)$, they
unlocked online dynamic minimum spanning tree maintenance in $\Oh(m \log n)$ time
(greedily cutting the heaviest edge on a cycle whenever a new edge is added),
and online connectivity queries for dynamic forests.

\textbf{Honest Downsides and Limitations:}
\begin{itemize}
  \item \emph{Amortized, not worst-case:} Splay-based link-cut trees achieve $\Oh(\log n)$
        amortized cost. While Sleator and Tarjan later developed a worst-case $\Oh(\log n)$
        variant using globally biased search trees, it is notoriously complicated, with
        formidable constant factors and code complexity that deter practical adoption.
  \item \emph{Subtrees are difficult:} Link-cut trees natively navigate root-to-node
        paths. Maintaining aggregates over \emph{subtrees} (such as the size or sum of
        an entire branch) requires tracking ``virtual subtree'' pointers that must be
        laboriously updated on every preferred-child splice.
  \item \emph{Forests only:} Link-cut trees cannot accommodate cycles. A single cycle
        causes \textsc{Access} to loop endlessly across path-parent edges.
\end{itemize}

\textbf{Where We Go Next (Lowest Common Ancestor and Range Structures):}
Throughout our study of tree decompositions -- from heavy-light chains to Euler tours and dynamic link-cut trees -- one primitive has appeared repeatedly: finding the \emph{Lowest Common Ancestor} ($\mathrm{LCA}$) of two nodes. While link-cut trees compute LCAs dynamically through successive accesses, static trees admit simpler solutions: binary lifting in $\Oh(\log n)$ query time, and the Bender--Farach-Colton reduction to $\pm 1$ Range Minimum Queries, answering LCA queries in $\Oh(1)$ worst-case time after linear $\Oh(n)$ preprocessing. In Chapter~23, we examine these static range structures and LCA algorithms, completing our repertoire for tree queries.
\end{connectionbox}

\section{Exercises}
\label{ch22:exercises}

\begin{exercise}[Quick mental check: splay key invariant]
What is the search key in an auxiliary splay tree? Can a node ever have a smaller search key than its ancestor in the represented tree?
\end{exercise}

\begin{exercise}[Quick mental check: path-parent vs parent]
In a link-cut tree node, what distinguishes a normal parent pointer from a path-parent pointer? How does a splay root detect that it is at the root of an auxiliary tree?
\end{exercise}

\begin{exercise}
Implement the splay-based link-cut tree with \textsc{MakeTree}, \textsc{Link},
\textsc{Cut}, \textsc{FindRoot}, \textsc{Evert}, and a path-sum query. Test it
against a brute-force forest (recompute paths by walking parent pointers) on
$10^4$ random operations.
\end{exercise}

\begin{exercise}
Prove that after $\textsc{Access}(v)$ the leftmost node of $v$'s auxiliary tree is
the represented root of $v$'s tree. Where in the proof do you use the
depth-ordering invariant?
\end{exercise}

\begin{exercise}
Trace $\textsc{Access}(9)$ on the post-$\textsc{Access}(4)$ forest drawn in
Section~\ref{ch22:access}. List every path-parent rewire and draw the final
preferred-edge set. How many preferred-child changes occurred?
\end{exercise}

\begin{exercise}
(Break it.) Implement \textsc{Evert} via the flip bit but \emph{forget} to push
the flip down inside \textsc{Splay}'s rotation. Construct a small sequence of
operations on which a path-sum query returns a wrong answer, and explain which
invariant was violated.
\end{exercise}

\begin{exercise}
Extend the structure to support a \emph{path range-add}: add a constant $c$ to
every node on the $u$--$v$ path. Specify the lazy tag, its \textsc{push} rule, and
how it interacts with the \texttt{flip} bit. Prove the aggregate stays correct.
\end{exercise}

\begin{exercise}
Show how to support \emph{subtree} aggregates (the sum over the whole represented
subtree of $v$) in a link-cut tree. (Hint: this is genuinely harder than path
aggregates; you must maintain a separate ``virtual subtree'' aggregate that the
\textsc{Access} splices keep up to date. Compare with Euler-tour trees, which
handle subtrees natively but not paths.)
\end{exercise}

\begin{exercise}
Give the LCA algorithm using two accesses precisely. Prove that the node returned
is the lowest common ancestor, and bound its cost. What goes wrong if the tree is
unrooted because you have been calling \textsc{Evert}?
\end{exercise}

\begin{exercise}
Use a link-cut tree to maintain the maximum edge weight on the path between two
nodes of a growing tree, and use it to verify in $\Oh(m\log n)$ whether a given
spanning tree is a minimum spanning tree of an $m$-edge graph. Describe the values
you store -- recall link-cut trees store \emph{node} values, so how do you encode
edge weights?
\end{exercise}

\begin{exercise}
(Worst case.) The splay-based structure is only $\Oh(\log n)$ \emph{amortized}.
Sketch how biased binary trees give an $\Oh(\log n)$ \emph{worst-case} bound per
operation, and explain why the amortized version is nonetheless preferred in
practice.
\end{exercise}

\begin{exercise}
Prove the telescoping step of Theorem~\ref{ch22:thm-amortized} carefully: state
precisely why $r(w)\ge r(t)$ when $w$ is the spliced-in parent of a path with
auxiliary root $t$, and where the ``$+1$ per splice'' units come from in the
Access Lemma.
\end{exercise}
